% ===============================================================
%  CHAPITRE 4 — LIMITES, PERSPECTIVES ET CONCLUSION (~1 page)
% ===============================================================
\chapter[Discussion et conclusion]{Limites, perspectives et conclusion}
\label{chap:discussion}

\section{Limites de l'étude}

La principale limite tient à la forme horaire de la consommation. Les deux méthodes que nous avons comparées reposent sur des profils régionaux, et leur accord montre seulement qu'elles sont cohérentes entre elles, pas qu'elles reproduisent la consommation réelle des Mauges. Les procédés industriels locaux, notamment agroalimentaires, peuvent avoir des rythmes propres. Seule une courbe de charge mesurée aux postes sources, que le gestionnaire de réseau pourrait fournir, permettrait de lever ce doute. La deuxième limite concerne le réseau. La charge des postes a été estimée par une répartition proportionnelle et confrontée à une capacité de planification, sans calcul de flux ni prise en compte des contraintes de tension. La troisième limite est économique. Seul le coût de la batterie a été chiffré, et non celui des centrales ni les émissions de CO\textsubscript{2} évitées, pourtant essentiels à une décision publique. Enfin, plusieurs hypothèses restent fragiles. Le parc est supposé constant sur les cinq années, l'horizon 2035 dépend d'hypothèses d'électrification empruntées à d'autres projets, et le nombre de ballons électriques, estimé faute de donnée locale, donne un potentiel technique qui ne préjuge pas de l'adhésion des ménages.

\section{Perspectives}

Ces limites dessinent les suites à donner à l'étude. La première serait de chiffrer le coût complet des trajectoires et le contenu carbone de l'électricité évitée, à partir du facteur d'émission horaire publié par éCO2mix. La deuxième serait d'engager avec Enedis et RTE une étude de raccordement sur les postes de Recouvrance, Vallet, Saint-Georges et Cholet. La troisième relève du dialogue avec les habitants, qu'il s'agisse de localiser les futurs parcs éoliens dans le cadre des zones d'accélération ou d'expérimenter le pilotage des ballons d'eau chaude avec quelques centaines de foyers volontaires. Le simulateur, dont toutes les hypothèses sont modifiables, peut accompagner ces échanges.

\section{Conclusion}

Ce projet devait déterminer jusqu'où Mauges Communauté peut développer le solaire et l'éolien sans perdre plus de 5~\% de sa production ni saturer son réseau. En reconstruisant 43~824 heures de consommation et de production et en confrontant chaque résultat à des ordres de grandeur indépendants, nous montrons que doubler le solaire ne suffit pas, car le solaire produit au même moment que le solaire déjà installé. La batterie demandée par la fiche règle le problème technique mais ne se rembourse jamais, quelles que soient les hypothèses de prix retenues. En misant d'abord sur l'éolien, qui produit quand le territoire consomme, et en pilotant les ballons d'eau chaude des ménages, la part de la consommation couverte heure par heure peut passer de 41,6 à 55,4~\% sans stockage et dans le respect de la limite d'écrêtement chaque année. Cette trajectoire n'est cependant réaliste qu'à deux conditions, orienter les projets vers les postes sources qui disposent encore de capacité et construire avec les habitants l'acceptation des nouveaux parcs éoliens. La question n'était donc pas seulement technique. Elle engage la répartition des installations entre les communes et la participation des ménages, que les élus disposent désormais d'un outil pour discuter.
